\subsection*{70. The Owner of the Customer Relationship Changes the Economics}

A D2C provider bears much of the cost of discovering, persuading, enrolling, and supporting each learner. A college may supply many learners at once, reducing direct customer acquisition expenditure but adding institutional procurement, curriculum coordination, faculty integration, reporting, and payment delays. An enterprise client may purchase a team engagement but demand customization, confidentiality, and implementation work.

The price per learner or participant is therefore not directly comparable across channels. In institutional contracts, the payer may purchase \textbf{reliability, reporting, and reduced administrative burden} alongside instruction. The business has not necessarily become more scalable merely because 1 contract covers many students.

\subsection*{71. Growth Can Require Adding Costs Before Adding Revenue}

Hiring a 2nd instructor, building an assessment bank, implementing access controls, or financing customer receivables may precede new sales. This creates a \textbf{step-investment problem}: a firm must spend to cross a threshold without knowing whether demand will fill the added capacity.

Use staged commitments where feasible: contract faculty by deliverable, pilot the rubric before building a platform, obtain deposits or signed scopes without misleading customers, and preserve the ability to stop if evidence disappoints. The relevant question is not simply expected profit after scaling, but the \textbf{maximum cash and reputation at risk before the next informative decision}.

\subsection*{72. Elastic Capacity Has Advantages: and Disadvantages}

Freelance reviewers, partner faculty, cloud infrastructure, and pay-per-use tools can make costs more flexible. They may reduce financial exposure when demand is volatile. But flexibility can mean higher unit prices, lower continuity, inconsistent standards, scheduling uncertainty, and dependence on suppliers outside the founder's control.

Use flexible capacity for genuinely variable work while investing in training, quality assurance, and stable relationships. A school whose entire teaching function appears only when enrollments are high may struggle to build faculty culture or reliable improvement routines. The goal is \textbf{appropriate risk sharing}, not avoidance of every fixed commitment forever.

\subsection*{73. Operating Leverage Can Be Intellectual Before It Is Financial}

An experienced founder can reuse a method of diagnosing projects in different settings. A trained faculty member can teach a 2nd course using the institute's standards. A robust evaluation harness can assess many different student artifacts. These may produce a rising ratio of \textbf{useful output to skilled labor} long before the company has large revenues.

Financial leverage appears when those capabilities translate into repeatable contribution and cash flow. This separation matters: building a better assessment system may initially cost money yet create a durable asset that improves future quality, speed, and credibility. In a mission-driven institution, some investments should be evaluated over a longer horizon than 1 cohort, while still respecting cash survival.

\subsection*{74. A Three-Dimensional View of Scale}

There are at least 3 ways to scale an institution:

\paragraph*{Throughput scale:} serve more learners, clients, cases, or units per period.

\paragraph*{Capability scale:} produce a broader or stronger range of verified outcomes with similar resources.

\paragraph*{Organizational scale:} enable other competent people and teams to deliver and improve the work without depending on the founder.

A video platform may increase throughput while doing little for capability. Better faculty training may increase organizational scale before learner count rises. A superior assessment-and-feedback loop may improve capability even if cohorts stay small. Good management chooses which dimension matters for the current stage and how the others support it.

\subsection*{75. The Transition from Founder Labor to Institutional Capital}

At the beginning, the founder's personal expertise compensates for missing systems. They explain concepts, notice errors, call customers, solve technical failures, and write proposals. Over time, the institution should transform recurring parts of that labor into \textbf{assets and trained capabilities}: explicit pedagogy, documented processes, quality assurance, technology, trusted relationships, and financial reserves.

The founder should not make themselves unnecessary overnight. Their role should move toward the highest-leverage problems: discovering new customer needs, designing the next capability, training senior faculty, assuring standards, and allocating capital. The institution becomes more durable when its ordinary promises no longer require extraordinary intervention by 1 individual.
